\documentclass[../../../include/open-logic-section]{subfiles}

\begin{document}
\olfileid{sth}{ord-arithmetic}{intro}

\olsection{جان پہچان}

\olref[ordinals][]{chap} وچ اسیں ترتیبی عدداں دا نظریہ بنایا سی۔
\olref[spine][]{chap} وچ اسیں ویکھیا کہ ترتیبی عدداں نوں ریڑھ دی ہڈی وانگ سمجھ سکدے آں، جیہدے دوالے درجہ وار سلسلے دا باقی حصہ بندا اے۔ پر ترتیبی عدداں دا کم بس ایہو نہیں۔ اوہناں اُتے حساب کرن دا کم وی اے۔

اسیں ایس ول پہلاں ای \olref[ordinals][idea]{sec} وچ اشارہ کیتا سی، جدوں اسیں $\omega$، $\omega+1$ تے $\omega+\omega$ دی گل کیتی سی۔ اوس ویلے اسیں غیر رسمی طور اُتے گل کیتی سی۔ ہُن ایس نوں ٹھیک طرح کھولھ کے دسّن دا ویلا آ گیا اے۔ پر ایہ دس دیئیے کہ ایس باب وچ فلسفہ بہت نہیں۔ ایتھے فنی گلّاں اگے ودھدیاں نیں، تے نال اک (کجھ) دلچسپ گل ویکھاں گے کہ اسیں اکّو کم دو وکھرے طریقیاں نال کر سکدے آں۔

\end{document}
